\documentclass[14pt]{extarticle}

\usepackage{graphicx}
\usepackage{indentfirst}
\usepackage[utf8]{inputenc}
\usepackage[russian]{babel}
\usepackage[left=3cm, right=2cm, top=2cm, bottom=2cm]{geometry}
\usepackage{setspace}
\usepackage{float}
\usepackage{caption}
\usepackage{amsmath}
\usepackage{enumitem}
\usepackage{multirow}
\usepackage{nicematrix}
\usepackage{makecell}
\usepackage{subcaption}
\usepackage{tikz}
\usepackage{titlesec}
\usepackage{tabularx}
\usepackage{array}
\usepackage{ragged2e}  % Для красивого выравнивания в колонках
\usepackage{adjustbox}
\usepackage{minted}
\usepackage{newfloat}
\usepackage{listings}
\usepackage{tcolorbox}

\tcbuselibrary{minted, breakable}
\usetikzlibrary{shapes.geometric, positioning, arrows.meta}
\renewcommand\tabularxcolumn[1]{m{#1}}

\floatname{listing}{Листинг}           % Название вместо "Listing"
\newfloat{listing}{htbp}{lop}          % Создаёт новый тип плавающего объект

% Настройка подписей для таблиц, рисунков и листингов через точку
\captionsetup[table]{labelsep=period}
\captionsetup[figure]{labelsep=period}
\captionsetup[listing]{labelsep=period}

% Общий стиль листинга
\lstset{
    basicstyle=\ttfamily\small,         % моноширинный шрифт
    keywordstyle=\color{blue}\bfseries, % ключевые слова синим и жирным
    stringstyle=\color{red},            % строки красным
    commentstyle=\color{green!50!black},% комментарии зеленым
    numbers=left,                        % номера строк слева
    numberstyle=\small\color{gray},       % стиль номеров
    stepnumber=1,                         % шаг нумерации
    numbersep=5pt,                        % расстояние между номерами и кодом
    showspaces=false,
    showstringspaces=false,
    frame=single,                         % рамка вокруг кода
    breaklines=true,                      % перенос длинных строк
    breakatwhitespace=true,               % перенос только на пробелах
    tabsize=4,                             % размер табуляции
    captionpos=b,                          % подпись под листингом
    escapeinside={@}{@}                   % вставка LaTeX-команд
}

\setminted{
    linenos,
    frame=single,
    framesep=2mm,
    breaklines,
    fontsize=\small,
    autogobble,
    breakanywhere
}
\begin{document}

\begin{center}

Министерство образования и науки Российской Федерации \\ \vspace{0.5cm}

Федеральное государственное автономное образовательное учреждение высшего образования \\ \vspace{1cm}

САНКТ-ПЕТЕРБУРГСКИЙ НАЦИОНАЛЬНЫЙ ИССЛЕДОВАТЕЛЬСКИЙ УНИВЕРСИТЕТ ИНФОРМАЦИОННЫХ ТЕХНОЛОГИЙ, МЕХАНИКИ И ОПТИКИ \\ \vspace{1cm}

Факультет программной инженерии и компьютерных технологий \\ \vspace{1cm}

\begin{spacing}{1.5}
Отчёт по лабораторной работе №1 \\
\end{spacing}

\end{center}

\vspace{3cm}

\begin{spacing}{1.5}
\begin{flushright}
\textbf{Выполнил:} Семикозов Ян Леонидович \\
\textbf{Группа:} P4119 \\
\textbf{Преподаватель:} Табунщик Сергей Михайлович \\
\end{flushright}
\end{spacing}

\thispagestyle{empty}
\clearpage
\setcounter{page}{1}

\section{Цель работы}
Получить общее представление о системе команд RISC-V путём решения несложной прикладной задачи -- найти наименьший элемент на диагонали в матрице $6\times6$.

\section{Ход работы}

В данной работе используется архитектура RISC-V с набором команд RV32I.
В соответствии с \textit{соглашением о вызовах функций RISC-V} аргументы функции передаются через регистры \texttt{a0-a7}, а результат возвращается в регистрах \texttt{a0, a1}.
Ниже приведена сводная таблица используемых регистров и их назначение в программе (таблица \ref{tab:regusage}).
\begin{table}[H]
\centering
\caption{Назначение используемых регистров}
\label{tab:regusage}
\small
\begin{tabularx}{\textwidth}{
|>{\centering\arraybackslash}X
|>{\centering\arraybackslash}X
|>{\centering\arraybackslash}X|}
\hline
\textbf{Регистры} & \textbf{Номер (x\#)} & \textbf{Назначение в коде} \\
\hline
\texttt{ra} & x1 & Адрес возврата из функции \\
\hline
\texttt{sp} & x2 & Указатель на вершину стека \\
\hline
\texttt{s0} & x8 & Адрес матрицы \\
\hline
\texttt{s1} & x9 & Текущее минимальное значение \\
\hline
\texttt{s2} & x10 & Индекс цикла i \\
\hline
\texttt{t0} & x5 & Граница цикла \\
\hline
\texttt{t1} & x6 & Смещение в матрице \\
\hline
\texttt{t2} & x7 & Текущий элемент в матрице \\
\hline
\texttt{a0} & x10 & Аргумент функции/код системного вызова, возвращаемое значение \\
\hline
\texttt{a1} & x11 & Аргумент для \texttt{ecall} \\
\hline
\end{tabularx}
\end{table}

\clearpage

Прежде чем описывать функциональную часть программы разместим входные и неизменяемые данные в секции \texttt{.data}.

\begin{listing}[H] 
    \begin{minted}[fontsize=\small, linenos, frame=single, breaklines]{asm}
	.data
    matrix:
        .word 12, 7, 31, 4, 5, 6
        .word 8, 10, 2, 11, 1, 9
        .word 14, 15, 6, 13, 12, 7
        .word 16, 17, 18, 5, 19, 20
        .word 21, 22, 23, 24, 9, 25
        .word 26, 27, 28, 29, 30, 10

    minelemstr:
        .asciiz "min element: "
    \end{minted} 
\caption{Код программы на языке ассемблера RISC-V} 
\label{lst:asmcode_dataseg} 
\end{listing}

Все инструкции программы разместим в секции \texttt{.text}; для глобальной видимости функций используется метка \texttt{.globl}.
В функции \textit{findMinElement} вначале создадим стековый фрейм для сохранения регистров, которые будут изменяться. 
При этом выделим 16 байт на стек и сохраним в него регистры \texttt{ra}, \texttt{s0}, \texttt{s1} и \texttt{s2}.
При завершении функции восстановим сохранённые значения регистров и освободим стековый фрейм.
Ниже в листинге \ref{lst:asmcode_findminelem} приведён полный код функции \textit{findMinElement}:
\begin{listing}[H] 
    \begin{minted}[fontsize=\small, linenos, frame=single, breaklines]{asm}
	.text
        .globl findMinElement
    findMinElement:
        addi sp, sp, -16
        sw ra, 12(sp)
        sw s0, 8(sp)
        sw s1, 4(sp)
        sw s2, 0(sp)
        mv s0, a0
        lw s1, 0(s0)
        li s2, 1

    loop:
        li t0, 6
        bge s2, t0, endloop
        li t1, 7
        mul t1, t1, s2
        slli t1, t1, 2
        add t1, s0, t1
        lw t2, 0(t1)
        blt t2, s1, updateMin
        j next_i

    updateMin:
        mv s1, t2

    next_i:
        addi s2, s2, 1
        j loop

    endloop:
        mv a0, s1
        lw ra, 12(sp)
        lw s0, 8(sp)
        lw s1, 4(sp)
        lw s2, 0(sp)
        addi sp, sp, 16
        ret
    \end{minted} 
\caption{Код функции \textit{findMinElement} на языке ассемблера RISC-V} 
\label{lst:asmcode_findminelem} 
\end{listing}

Ниже в листинге \ref{lst:asmcode_main} приведён полный код функции \textit{main}, которая вызывает функцию поиска минимального элемента и выводит результат на экран с использованием системных вызовов.
Код системного вызова записывается в регистр \texttt{a7}, a аргумент в \texttt{a1}:
\begin{itemize}
    \item Код 1 -- вывод целого числа.
    \item Код 4 -- вывод строки.
    \item Код 11 -- вывод символа.
\end{itemize}

\begin{listing}[H] 
    \begin{minted}[fontsize=\small, linenos, frame=single, breaklines]{asm}
    main:
        addi sp, sp, -16
        sw ra, 12(sp)

        la a0, matrix
        call findMinElement

        mv t0, a0
    
        # print string
        la a1, minelemstr
        li a0, 4
        ecall

        # print integer
        mv a1, t0
        li a0, 1
        ecall

        # print newline
        li a1, 10
        li a0, 11
        ecall


        lw ra, 12(sp)
        addi sp, sp, 16
        li a0, 0
        ret

    \end{minted} 
\caption{Код функции \textit{main} на языке ассемблера RISC-V} 
\label{lst:asmcode_main} 
\end{listing}

Ниже в листинге \ref{lst:ccode} приведён полный код программы на языке C, решающей ту же задачу.
\begin{listing}[H]
    \begin{minted}[fontsize=\small, linenos, frame=single, breaklines]{c}
    #include <stdio.h>

    #define N 6
    #define M 6

    int findMinElement(int matrix[N][M])
    {
        int min = matrix[0][0];

        for (int i = 1; i < N && i < M; i++) {
            if (matrix[i][i] < min) {
                min = matrix[i][i];
            }
        }

        return min;
    }

    int main()
    {
        int matrix[N][M] = {
            { 12, 7, 31, 4, 5, 6 },
            { 8, 10, 2, 11, 1, 9 },
            { 14, 15, 6, 13, 12, 7 },
            { 16, 17, 18, 5, 19, 20 },
            { 21, 22, 23, 24, 3, 25 },
            { 26, 27, 28, 29, 30, 4 }
        };

        int min = findMinElement(matrix);

        printf("min element: %d\n", min);

        return 0;
    }
    \end{minted}
\caption{Код программы на C}
\label{lst:ccode}
\end{listing}

Сборка и запуск программы осуществлялись с использованием собранного кросс-компилятора \textit{RISC-V} и эмулятора \textit{QEMU}. 
Ниже приведены команды, использованные для сборки и запуска программы из корня проекта:

\begin{center}
\begin{minipage}{0.7\linewidth}
\begin{Verbatim}[frame=single]
make -C src && qemu-riscv32 src/main.elf
\end{Verbatim}
\end{minipage}
\end{center}

Переписанный код на языке ассемблера был проверен в симуляторе \textit{Venus}.

\section{Выводы}

В ходе выполнения лабораторной работы изучены базовые принципы написания кода под архитектуру RISC-V. 
Реализована программа на языке C и её эквивалент на языке ассемблера. 
Ознакомились с процессом сборки и запуска программ под архитектуру RISC-V. 

\end{document}
